% TODO make hw stylesheet
\documentclass[11pt]{scrartcl}
\usepackage[a4paper, total={6in, 8in}]{geometry}
\usepackage[usenames,dvipsnames,svgnames]{xcolor}
\usepackage[shortlabels]{enumitem}
\usepackage[framemethod=TikZ]{mdframed}
\usepackage{amsmath,amssymb,amsthm}
\usepackage[colorlinks]{hyperref}
\usepackage{multicol}
\usepackage{tabularx}

\hypersetup{
	colorlinks=true,
	linkcolor=blue,
	filecolor=magenta,      
	urlcolor=magenta,
	pdftitle={MAT 322 HW}
}

\usepackage{microtype}
\usepackage{mathtools}
\usepackage[headsepline]{scrlayer-scrpage}
\usepackage{thmtools}
\usepackage[textsize=footnotesize]{todonotes}

\mdfdefinestyle{mdbluebox}{roundcorner=10pt,innerbottommargin=9pt,
	linecolor=blue,backgroundcolor=TealBlue!5,}
\declaretheoremstyle[headfont=\sffamily\bfseries\color{MidnightBlue},
mdframed={style=mdbluebox},]{thmbluebox}

\mdfdefinestyle{mdbloobox}{frametitlefont=\bfseries,innerbottommargin=8pt,
	nobreak=true,backgroundcolor=TealBlue!5,linecolor=blue,}
\declaretheoremstyle[headfont=\bfseries\color{MidnightBlue},
mdframed={style=mdbloobox},headpunct={\\[3pt]},postheadspace=0pt,]{thmbloobox}

\mdfdefinestyle{mdredbox}{frametitlefont=\bfseries,innerbottommargin=8pt,
	nobreak=true,backgroundcolor=Salmon!5,linecolor=RawSienna,}
\declaretheoremstyle[headfont=\bfseries\color{RawSienna},
mdframed={style=mdredbox},headpunct={\\[3pt]},postheadspace=0pt,]{thmredbox}

\mdfdefinestyle{mdgreenbox}{linecolor=ForestGreen,backgroundcolor=ForestGreen!5,
	linewidth=2pt,rightline=false,leftline=true,topline=false,bottomline=false,}
\declaretheoremstyle[headfont=\bfseries\sffamily\color{ForestGreen!70!black},
mdframed={style=mdgreenbox},headpunct={ --- },]{thmgreenbox}

\mdfdefinestyle{mdorangebox}{linecolor=RedOrange,backgroundcolor=RedOrange!5,
	linewidth=2pt,rightline=false,leftline=true,topline=false,bottomline=false,}
\declaretheoremstyle[headfont=\bfseries\sffamily\color{RedOrange!70!black},
mdframed={style=mdorangebox},headpunct={ --- },]{thmorangebox}

\mdfdefinestyle{mdblackbox}{linecolor=black,backgroundcolor=RedViolet!5!gray!5,
	linewidth=3pt,rightline=false,leftline=true,topline=false,bottomline=false,}
\declaretheoremstyle[mdframed={style=mdblackbox}]{thmblackbox}

\declaretheoremstyle[spaceabove=6pt,spacebelow=6pt]{basehead}

\declaretheorem[style=basehead,name=Problem]{problem}
\declaretheorem[style=thmredbox,name=Problem,numbered=no]{problem*}
\declaretheorem[style=thmbloobox,name=Setup]{setup} % for config geo
\declaretheorem[style=thmredbox,name=Example,sibling=problem]{example}
\declaretheorem[style=thmredbox,name=$\bigstar$ Problem,sibling=problem]{niceprob}
\declaretheorem[style=thmbluebox,name=Theorem,sibling=problem]{theorem}
\declaretheorem[style=thmbluebox,name=Corollary,sibling=theorem]{corollary}
\declaretheorem[style=thmbluebox,name=Proposition,sibling=theorem]{prop}
\declaretheorem[style=thmbluebox,name=Lemma,numberwithin=problem]{lemma}
\declaretheorem[style=thmbluebox,name=Theorem,numbered=no]{theorem*}
\declaretheorem[style=thmbluebox,name=Lemma,numbered=no]{lemma*}
\declaretheorem[style=thmgreenbox,name=Claim,numberwithin=problem]{claim}
\declaretheorem[style=thmgreenbox,name=Claim,numbered=no]{claim*}
\declaretheorem[style=thmblackbox,name=Remark,numberwithin=problem]{remark}
\declaretheorem[style=thmblackbox,name=Remark,numbered=no]{remark*}
\declaretheorem[style=thmgreenbox,name=Definition,sibling=theorem]{definition}
\declaretheorem[style=thmgreenbox,name=Definition,numbered=no]{definition*}
\declaretheorem[style=thmorangebox,name=Warning,sibling=theorem]{warning}
\declaretheorem[style=thmorangebox,name=Warning,numbered=no]{warning*}
\declaretheorem[style=thmorangebox,name=Note,numbered=no]{note}
\declaretheorem[style=thmblackbox,name=Exercise,sibling=problem]{exercise}
\declaretheorem[style=thmblackbox,name=Exercise,numbered=no]{exercise*}
\declaretheorem[style=thmblackbox,name=Question,sibling=problem]{question}
\declaretheorem[style=thmblackbox,name=Question,numbered=no]{question*}

\addtolength{\textheight}{3.14cm}
\providecommand{\ol}{\overline}
\providecommand{\ul}{\underline}
\providecommand{\eps}{\varepsilon}
\providecommand{\half}{\frac{1}{2}}
\providecommand{\dang}{\measuredangle} %% Directed angle
\providecommand{\CC}{\mathbb C}
\providecommand{\FF}{\mathbb F}
\providecommand{\NN}{\mathbb N}
\providecommand{\QQ}{\mathbb Q}
\providecommand{\RR}{\mathbb R}
\providecommand{\ZZ}{\mathbb Z}
\providecommand{\dg}{^\circ}
\providecommand{\ii}{\item}
\providecommand{\setdiff}{\smallsetminus}
\providecommand{\alert}[1]{{\sffamily\textbf{\textcolor{blue}{#1}}}}
\providecommand{\inv}{^{-1}}
\providecommand{\mb}[1]{\mathbf{#1}}
\providecommand{\dif}{\;d}

\reversemarginpar

\newenvironment{soln}{\begin{proof}[Solution]}{\end{proof}}
\newenvironment{walkthrough}{\noindent \textbf{Walkthrough.}}{\medskip}
\newlist{walk}{enumerate}{3}
\setlist[walk]{label=\bfseries (\alph*)}

\providecommand{\pow}{\operatorname{Pow}}
\providecommand{\vocab}[1]{{\textbf{\textcolor{ForestGreen}{#1}}}}
\providecommand{\lcm}{\operatorname{lcm}}
\providecommand{\abs}[1]{\left|#1\right|}

\usepackage{asymptote}
\begin{asydef}
	defaultpen(fontsize(10pt));
	unitsize(1cm);
	size(8cm); // set a reasonable default
	usepackage("amsmath");
	usepackage("amssymb");
	settings.tex="pdflatex";
	settings.outformat="pdf";
	// Replacement for olympiad+cse5 which is not standard
	import geometry;
	// recalibrate fill and filldraw for conics
	void filldraw(picture pic = currentpicture, conic g, pen fillpen=defaultpen, pen drawpen=defaultpen)
	{ filldraw(pic, (path) g, fillpen, drawpen); }
	void fill(picture pic = currentpicture, conic g, pen p=defaultpen)
	{ filldraw(pic, (path) g, p); }
	// some geometry
	pair foot(pair P, pair A, pair B) { return foot(triangle(A,B,P).VC); }
	pair orthocenter(pair A, pair B, pair C) { return orthocentercenter(A,B,C); }
	pair centroid(pair A, pair B, pair C) { return (A+B+C)/3; }
	// cse5 abbreviations
	path CP(pair P, pair A) { return circle(P, abs(A-P)); }
	path CR(pair P, real r) { return circle(P, r); }
	pair IP(path p, path q) { return intersectionpoints(p,q)[0]; }
	pair OP(path p, path q) { return intersectionpoints(p,q)[1]; }
	path Line(pair A, pair B, real a=0.6, real b=a) { return (a*(A-B)+A)--(b*(B-A)+B); }
	// cse5 more useful functions
	picture CC() {
		picture p=rotate(0)*currentpicture;
		currentpicture.erase();
		return p;
	}
	pair MP(Label s, pair A, pair B = plain.S, pen p = defaultpen) {
		Label L = s;
		L.s = "$"+s.s+"$";
		label(L, A, B, p);
		return A;
	}
	pair Drawing(Label s = "", pair A, pair B = plain.S, pen p = defaultpen) {
		dot(MP(s, A, B, p), p);
		return A;
	}
	path Drawing(path g, pen p = defaultpen, arrowbar ar = None) {
		draw(g, p, ar);
		return g;
	}
\end{asydef}

\usepackage{mathrsfs}

\ihead{\footnotesize MAT 322 HW02}
\ohead{\footnotesize Last compiled \today}

\title{\vspace{-2em}MAT 322 HW02}
\author{\large Aditya Pahuja}
\date{\large Due 02/11}
\begin{document}
\maketitle

\begin{problem}
	Let $f\colon\RR^n\to\RR$ and $g\colon\RR^n\to\RR$ be functions.
	If $f$ and $g$ are differentiable
	(respectively, class $C^1$ on a neighborhood of $x$)
	at the point $x\in\RR^n$, prove that their sum $f+g$ and product $fg$
	are both differentiable
	(respectively, class $C^1$ on a neighborhood of $x$)
	at $x$ and find their derivatives at $x$.
\end{problem}

\begin{soln}
	If $f$ and $g$ are differentiable at $x$, then
	\begin{align*}
		\Delta_f(x,h) := f(x+h) - f(x) &= (\mathrm df)_x(h) + \|h\|\eps_f(h)\\
		\Delta_g(x,h) := g(x+h) - g(x) &= (\mathrm dg)_x(h) + \|h\|\eps_g(h)
	\end{align*}
	for functions $\eps_f,\eps_g\colon\RR^n\to\RR$ whose limits
	as $h$ approaches zero are both $0$.
	Then, we can see that the sum $f+g$ is differentiable at $x$
	with derivative $(\mathrm d(f+g))_x =
	\boxed{(\mathrm df)_x + (\mathrm dg)_x}$
	because
	\[(f+g)(x+h) - (f+g)(x) =
	(\mathrm df)_x(h) + (\mathrm dg)_x(h) +
	\|h\|(\eps_f(h) + \eps_g(h))\]
	and $\eps_f + \eps_g$ is zero in the limit as $h$ goes to zero.
	Similarly, the product $fg$ is differentiable with derivative
	$(\mathrm d(fg))_x =
	\boxed{f(x)\cdot(\mathrm dg)_x + g(x)\cdot(\mathrm df)_x}$.
	To show this, we write
	\begin{align*}
		(fg)(x+h) - (fg)(x) &= (fg)(x+h) -
		f(x+h)g(x) + f(x+h)g(x) - (fg)(x)\\
		&= f(x+h)\Delta_g(x,h) + g(x)\Delta_f(x,h)
	\end{align*}
	Now, we handle these two addends separately.
	For the first addend, we can expand
	$f(x+h) = f(x) + (\mathrm df)_x(h) + \|h\|\eps_f(h)$
	and $\Delta_g(x,h) = (\mathrm dg)_x(h) + \|h\|\eps_g(h)$.
	This leaves four important pieces:
	\begin{itemize}
		\ii $f(x)\cdot(\mathrm dg)_x(h)$;
		\ii $(\mathrm df)_x(h)\cdot(\mathrm dg)_x(h)$;
		\ii $\|h\|\eps_f(h)\cdot (\mathrm dg)_x(h)$; and
		\ii $\|h\|\eps_g(h)\cdot f(x+h)$.
	\end{itemize}
	Of these, I claim that the latter three each decay faster than $h$
	as $h$ goes to zero, i.e. they are part of the error term
	when computing the derivative of $fg$. We verify this:
	\begin{itemize}
		\ii Noting that $(\mathrm df)_x$ and $(\mathrm dg)_x$ are linear maps,
		they are continuous on all of $\RR^n$. To verify that
		\[\lim_{h\to 0} \frac{(\mathrm df)_x(h)}{\|h\|}\cdot(\mathrm dg)_x(h)
		= 0,\]
		we need only verify that
		\[\frac{(\mathrm df)_x(h)}{\|h\|} =
		(\mathrm df)_x\left(\frac{h}{\|h\|}\right)\]
		is bounded on a neighborhood of $0$ because $(\mathrm dg)_x(0) = 0$,
		and this follows from the continuity of linear functions.
		\ii The other two terms are easier: they already come with a $\|h\|$
		attached, so we only need to check that
		\[\lim_{h\to 0} \eps_f(h)\cdot(\mathrm dg)_x(h) =
		\lim_{h\to 0}\eps_g(h)\cdot f(x+h) = 0.\]
		Since $\eps_f(h)$ and $\eps_g(h)$ already approach $0$ with $h$,
		it suffices to check that $(\mathrm dg)_x(h)$
		and $f(x+h)$ are bounded on a neighborhood of $0$,
		but again this follows from continuity of the functions
		$(\mathrm dg)_x$ and $f$ (the latter is continuous at $x$
		because it is differentiable at $x$).
	\end{itemize}
	We will denote the sum of these three terms with $\|h\|\eps_1(h)$,
	keeping in mind that
	\[\lim_{h\to 0}\eps_1(h) = 0\]

	The other addend is much simpler to work with: it expands as
	\[g(x)\cdot(\mathrm df)_x(h) + \|h\|\eps_f(h)\cdot g(x).\]
	Noting that $g(x)$ is constant with respect to $h$, we directly have
	\[\lim_{h\to 0} \eps_f(h)\cdot g(x) = g(x)\lim_{h\to0}\eps_f(h) = 0.\]
	We will write $\eps_f(h)g(x) = \eps_2(h)$, noting also that
	$\eps_2(h)$ approaches $0$ as $h$ approaches $0$.
	
	Thus, we can conclude that
	\[(fg)(x+h) - (fg)(x) =
	f(x)\cdot(\mathrm dg)_x(h) + g(x)\cdot (\mathrm df)_x(h)
	+ \|h\| (\eps_1(h) + \eps_2(h)),\]
	which means that, by the definition of derivative,
	the derivative of $fg$ at $x$ is
	$f(x)\cdot(\mathrm dg)_x + g(x)\cdot(\mathrm df)_x$.
	
	In the case that $f$ and $g$ are both class $C^1$
	on a neighborhood of $x$, we can use
	the fact that class $C^1$ functions are differentiable on this neighborhood
	together with the new derivative formulas we have just proven above
	to show that $f+g$ and $fg$ are also class $C^1$. Specifically,
	the Jacobian matrices of $(\mathrm d(f+g))_x$ and $(\mathrm d(fg))_x$
	precisely contain the partial derivatives
	of their respective functions, so we just need to verify
	that these entries are all continuous; but this is obvious for sums because
	\[(\mathrm d(f+g))_x = (\mathrm df)_x + (\mathrm dg)_x\]
	and so the relevant partials (matrix entries)
	are sums of continuous functions
	and it is obvious for products because
	\[(\mathrm d(fg))_x = f(x)\cdot (\mathrm dg)_x + g(x)\cdot(\mathrm df)_x,\]
	where the partials look like
	\[f(x)\cdot\frac{\partial g}{\partial x_j} +
	g(x)\cdot \frac{\partial f}{\partial x_j},\]
	which are continuous because they are sums of products
	of continuous functions.
\end{soln}

\pagebreak
\begin{problem}
	Show that the function $f(x,y) = |xy|$ is differentiable at $(0,0)$,
	but it is not of class $C^1$ on any open neighborhood of $(0,0)$.
\end{problem}

\begin{soln}
	I claim that $(\mathrm df)_{(0,0)}$ is the zero map.
	To check this we need to show that
	\[\lim_{(\alpha,\beta)\to (0,0)}
	\frac{f(\alpha,\beta) - f(0,0) - (\mathrm df)_{(0,0)}(\alpha,\beta)}
	{\|(\alpha,\beta)\|} =
	\lim_{(\alpha,\beta)\to(0,0)} = \frac{|\alpha\beta|}
	{\sqrt{\alpha^2+\beta^2}} = 0.\]
	By definition, this means we want to show that for each $\eps > 0$,
	there exists a $\delta > 0$ such that $\|(\alpha,\beta)-(0,0)\|=
	\sqrt{\alpha^2+\beta^2}<\delta$
	implies $\frac{|\alpha\beta|}{\sqrt{\alpha^2+\beta^2}} < \eps$.
	By the AM-GM inequality,
	\[\alpha^2 + \beta^2 \ge 2|\alpha\beta|\]
	for any reals $\alpha$, $\beta$, so, taking $\delta = 2\eps$, we get
	\[\sqrt{a^2+b^2}<\delta \implies \frac{|\alpha\beta|}{\sqrt{a^2+b^2}}
	\le \frac{\half(\alpha^2+\beta^2)}{\sqrt{\alpha^2+\beta^2}}
	= \frac{\sqrt{\alpha^2+\beta^2}}{2} < \half\delta = \eps.\]
	Therefore, the limit is $0$, so $f$ is differentiable at $(0,0)$.
	
	However, $f$ does not have continuous partials on a neighborhood
	of $(0,0)$, since, for example, $\frac{\partial f}{\partial x}(0,y)$
	doesn't even exist for $y\ne 0$.
\end{soln}

\begin{problem}
	Define $f\colon\RR\to\RR$ by setting $f(0) = 0$ and
	\[f(t) = t^2\sin(1/t)\text{ if }t\ne 0.\]
	\begin{enumerate}[(a)]
		\ii Show that $f$ is differentiable at $0$, and calculate $f'(0)$.
		\ii Calculate $f'(t)$ if $t\ne 0$.
		\ii Show $f'$ is not continuous at $0$.
		\ii Conclude that that $f$ is differentiable but
		not continuously differentiable on $\RR$.
	\end{enumerate}
\end{problem}

\begin{soln}
	The derivative of $f$ at $0$ is
	\[\lim_{t\to 0} \frac{f(t) - f(0)}{t-0}
	= \lim_{t\to 0} t\sin(1/t) = 0,\]
	where the second equality uses the fact that $\sin(1/t)$ is bounded
	on a neighborhood of $0$, since $|\sin(x)| \le 1$ for all $x\in\RR$.
	When $t\ne 0$, the derivative can be computed with
	the standard derivative laws to get
	\[f'(t) = t^2\cos(1/t)\cdot\frac{-1}{t^2} + 2t\sin(1/t)
	= 2t\sin(1/t) - \cos(1/t).\]
	As $t$ approaches $0$, we have
	\[\lim_{t\to0} f'(t) = \lim_{t\to0} 2t\sin(1/t) - \cos(1/t).\]
	When $t$ is small,
	$2t\sin(1/t)$ will become arbitrarily small
	(it is bounded in absolute value by $|2t|$). Thus, the limit equals
	\[-\lim_{t\to0} \cos(1/t) = -\lim_{y\to\infty} \cos(y),\]
	but this doesn't exist, so $f'$ cannot be continuous at $0$.
	Therefore, $f$ is differentiable, but its derivative isn't continuous,
	so it is not class $C^1$.
\end{soln}

\begin{problem}
	Let $f\colon U\to\RR^m$ be a function defined on an open set
	$U\subseteq\RR^m$, such that all partial derivatives
	$\partial f/\partial x_i$ exist and are bounded in a neighborhood
	of a point $x\in U$. Prove that $f$ is continuous at $x$.
\end{problem}

\begin{soln}
	To start with, we reduce to the $m=1$ case via the following claim:
	\begin{claim}
		The function $f\colon\RR^n\to\RR^m$ is continuous at a point $x$
		if each component function $f_i\colon\RR^n\to\RR$ is continuous at $x$.
	\end{claim}
	\begin{proof}
		If $f_1,\dots,f_m\colon\RR^n\to\RR$
		are all continuous at $x$, then the linear combination
		\[f = f_1e_1 + f_2e_2 + \cdots + f_me_m \colon\RR^n\to\RR^m\]
		is continuous at $x$ because it is a sum of continuous functions at $x$.
	\end{proof}

	Thus, it suffices to solve the problem for functions $f\colon\RR^n\to\RR$
	and then use the claim to patch together $m$ different
	$\RR^n\to\RR$ functions for the general setting.
	We will show in that
	\[\lim_{h\to 0} f(x+h) - f(x) = 0.\]
	Define
	\[v_i = h_1e_1 + h_2e_2 + \cdots + h_iv_i\]
	for $i = 1,2,\dots,n$, so $h = v_n$, and let $v_0=0$. Then, write
	\[f(x + h) - f(x) = \sum_{i=1}^nf(x + v_i) - f(x + v_{i-1}).\]
	Since each of the partials of $f$ exists on a neighborhood $N$ of $x$,
	we can use the mean value theorem
	to find a value of $t_i$ for $i=1,2,\dots,n$ such that
	\[f(x + v_i) - f(x + v_{i-1}) =
	\frac{\partial f}{\partial x_i}(x + v_{i-1} + t_ie_i)\cdot h_i\]
	where $0 < t < h_i$ as long as each $v_i$ lies in $N$
	(which just requires $\|x-h\|$ to be sufficiently small).
	Since all the partials are bounded on this neighborhood,
	we can find a constant $C$ such that, for each $i$,
	\[\abs{\frac{\partial f}{\partial x_i}(x + v_{i-1} + t_ie_i)} < C.\]
	Then, we simply have
	\[\lim_{h\to 0} f(x+h) - f(x) =
	\lim_{h\to 0}\sum_{i=1}^nf(x + v_i) - f(x + v_{i-1}) <
	\lim_{h\to 0} C\cdot(h_1 + \cdots + h_n) = 0,\]
	hence $f$ is continuous at $x$ if its partials
	exist and are bounded on a neighborhood of $x$.
\end{soln}

\pagebreak
\begin{problem}
	Suppose $f\colon U\to\RR^m$ and $g\colon V\to\RR^\ell$
	are class $C^1$ functions defined on open sets
	$U\subseteq\RR^n$, $V\subseteq\RR^m$ such that $f(U)\subseteq V$.
	Prove that the composition $g\circ f\colon U\to\RR^\ell$ is also $C^1$.
\end{problem}

\begin{soln}
	Since $f$ and $g$ are class $C^1$, they are differentiable,
	so their composition is also differentiable,
	and we can use the chain rule to compute that
	the derivative of $g\circ f$ at a point $x$ is
	\[(\mathrm d(g\circ f))_x = (\mathrm dg)_{f(x)}\cdot(\mathrm df)_x.\]
	The entries of the Jacobian matrices of
	$(\mathrm dg)_{f(x)}$ and $(\mathrm df)_x$
	are all continuous functions in $x$
	since $f$ is differentiable-hence-continuous
	and $f$, $g$ are class $C^1$, so the product of the matrices
	has entries that are also continuous in $x$
	because these entries are sums of products of continuous functions.
\end{soln}

\begin{problem}
	Let $f\colon\RR^3\to\RR^2$ and $g\colon\RR^2\to\RR^3$ be the functions
	\begin{align*}
		f(x) &= (e^{2x_1 + x_2}, 3x_2 - \cos x_1, x_1^2 + x_2 + 2)\\
		g(y) &= (3y_1 + 2y_2 + y_3^2, y_1^2 - y_3 + 1).
	\end{align*}
	\begin{enumerate}[(a)]
		\ii Compute $(\mathrm d(g\circ f))_0$ without computing $(g\circ f)(x)$.
		\ii Compute $(\mathrm d(f\circ g))_0$ without computing $(f\circ g)(x)$.
	\end{enumerate}
\end{problem}

\begin{soln}
	We have $f(0,0) = (1, -1, 2)$ and $g(0,0,0) = (0, 1)$.
	In general,
	\[(\mathrm df)_x = \begin{pmatrix}
		2e^{2x_1+x_2} & e^{2x_1 + x_2}\\
		\sin x_1 & 3\\
		2x_1 & 1
	\end{pmatrix},\qquad
	(\mathrm dg)_x = \begin{pmatrix}
		3 & 2 & 2y_3\\
		2y_1 & 0 & -1
	\end{pmatrix},\]
	so we can compute the following four derivatives:
	\begin{align*}
		(\mathrm df)_0 &= \begin{pmatrix}
			2&1\\0&3\\0&1
		\end{pmatrix}\\
		(\mathrm df)_{g(0)} &= \begin{pmatrix}
			2e&e\\0&3\\0&1
		\end{pmatrix}\\
		(\mathrm dg)_0 &= \begin{pmatrix}
			3&2&0\\0&0&-1
		\end{pmatrix}\\
		(\mathrm dg)_{f(0)} &= \begin{pmatrix}
			3&2&4\\2&0&-1
		\end{pmatrix}
	\end{align*}
	Therefore, the answer to (a) is
	\[(\mathrm d(g\circ f)) = (\mathrm dg)_{f(0)}\cdot(\mathrm df)_0
	= \begin{pmatrix}
		3&2&4\\2&0&-1
	\end{pmatrix}\begin{pmatrix}
		2&1\\0&3\\0&1
	\end{pmatrix} = \boxed{\begin{pmatrix}
		6&13\\4&1
	\end{pmatrix}}
	\]
	and the answer to (b) is
	\[(\mathrm d(f\circ g)) = (\mathrm df)_{g(0)}\cdot(\mathrm dg)_0
	= \begin{pmatrix}
		2e&e\\0&3\\0&1
	\end{pmatrix}\begin{pmatrix}
		3&2&0\\0&0&-1
	\end{pmatrix} = \boxed{\begin{pmatrix}
			6e&4e&-e\\0&0&-3\\0&0&-1
	\end{pmatrix}}.\qedhere
	\]
\end{soln}

\begin{problem}
	Let $f\colon\RR^n\to\RR^m$ be differentiable at the point $u\in\RR^n$,
	and $b\in\RR^m$ a vector. Consider the function $F\colon\RR^m\to\RR$
	given by
	\[F(x) = \|f(x) - b\|^2.\]
	Prove that $F$ is differentiable at $u$ and that, for any $v\in\RR^n$,
	\[(\mathrm dF)_u(v) = 2\langle f(u) - b, (\mathrm df)_u(v)\rangle.\]
\end{problem}

\begin{soln}
	Let $g(x) = \|x\|^2 = \langle x,x\rangle$ be a function from
	$\RR^m$ to $\RR$.
	\begin{claim}
		The function $g$ is differentiable,
		and the derivative of $g$ at point $u$ is
		the linear map determined by
		\[(\mathrm dg)_u(v) = \langle 2u,v\rangle\]
		for all $v\in\RR^m$.
	\end{claim}
	\begin{soln}
		This is just a simple computation:
		\begin{align*}
			\lim_{h\to 0} \frac{g(u+h)-g(u) - (\mathrm dg)_u(h)}{\|h\|}
			&= \lim_{h\to 0} \frac{\langle 2u+h, h\rangle - \langle 2u,h\rangle}
			{\|h\|}\\
			&= \lim_{h\to 0} \frac{\langle h,h\rangle}{\|h\|}
			= \lim_{h\to 0} \frac{\|h\|^2}{\|h\|}
			= \lim_{h\to 0} \|h\| = 0.
		\end{align*}
		Note that the matrix for this linear map is $2u^\top$, the transpose of
		the vector $2u$.
	\end{soln}
	
	Now we note that $f(x)-b$ is differentiable at $u$
	with derivative $(\mathrm df)_u$ because the constant function
	sending every point of $\RR^n$ to $b$ has a derivative of $0$.
	Therefore, the chain rule says that $F(x) = g(f(x) - b)$,
	being a composition of functions that are differentiable at $u$,
	is also differentiable at $u$, with derivative
	\[(\mathrm dF)_u = (\mathrm dg)_{f(u) - b}\cdot(\mathrm df)_u.\]
	Then, we conclude that
	\[(\mathrm dF)_u(v) = (\mathrm dg)_{f(u) - b}\cdot (\mathrm df)_u(v)
	= 2(f(u) - b)^\top \cdot(\mathrm df)_{u}(v)
	= 2\langle f(u) - b, (\mathrm df)_u(v)\rangle,\]
	as desired.
\end{soln}





















\end{document}